\documentclass[11pt]{article}

\usepackage[margin=1in]{geometry}
\usepackage{amsmath}
\usepackage{amssymb}
\usepackage{booktabs}
\usepackage{graphicx}
\usepackage[hidelinks]{hyperref}
\usepackage{microtype}
\usepackage{siunitx}

\sisetup{detect-all}

\title{Internal Fidelity and Real-Space Locality of\\Reduced Wannier Hamiltonians:\\A Rank-10 Silicon Case Study}
\author{Project Koios}
\date{Development technical report, version 2}

\newcommand{\Hsrc}{H_{\mathrm W}}
\newcommand{\Hred}{H_{\mathrm r}}
\newcommand{\kvec}{\mathbf{k}}
\newcommand{\Rvec}{\mathbf{R}}
\newcommand{\qvec}{\mathbf{q}}
\newcommand{\resultfigure}[4]{%
  \begin{figure}[p]
    \centering
    \includegraphics[width=\textwidth,height=0.78\textheight,keepaspectratio]{#1}
    \caption{#2 #3}
    \label{#4}
  \end{figure}%
}

\begin{document}
\maketitle

\begin{abstract}
We study how real-space reduction changes a fixed, rank-10 Wannier Hamiltonian derived
from a scalar-relativistic silicon calculation.  The purpose is not to validate the
Hamiltonian as a quantitatively accurate model of silicon.  Instead, it is to determine
which internal fidelity properties survive changes of gauge, literal removal of complete
Wigner--Seitz shells, and bounded refitting of retained shells.

A constructed momentum-dependent unitary leaves the complete represented spectrum
invariant to \(3.20\times10^{-14}\)~eV while changing restricted alignment residuals and
Fourier-block locality.  Literal nested shell deletion produces an almost
one-dimensional spectral--operator loss family and requires 371 of 383 representatives
for the tightest proper reduction.  Bounded shell refitting creates genuine
same-support spectral--operator tradeoffs and reduces the observed support at several
post-hoc error pairs.  However, errors measured on the 408-point fitting inventory
understate exact integrated operator error and independent dense-mesh spectral error.
Band-edge quantities are still more demanding: moderate supports substantially alter
conduction masses, split the nearly degenerate valence manifold, and amplify differences
among symmetry-equivalent high-symmetry points.

The results show that spectral agreement, operator agreement, real-space locality, and
band-edge fidelity are distinct properties.  Near-full support is required to preserve
all of them for this represented source.  Because the source itself fails separate
DFT-parent gap and electron-mass targets, every conclusion is restricted to internal
Wannier-model fidelity.
\end{abstract}

\paragraph{Evidence status.}
All numerical results in this report are checksum-bound, producer-checked development
outputs.  They have not been independently reconstructed.  No Quantum ESPRESSO or
Wannier90 process was invoked for the analyses reported here, and no frozen 402-point
confirmation values were consumed.

\section{Introduction}

Wannier Hamiltonians provide a finite-dimensional representation in which a sampled
electronic structure can be interpolated and expressed through real-space matrix
blocks~\cite{MarzariVanderbilt1997,SouzaMarzariVanderbilt2001,Mostofi2014,Yates2007}.
This representation is attractive for reduced models because the Hamiltonian can be
written as a Fourier series and approximated by a smaller set of real-space terms.
The apparent simplicity of that procedure hides several independent questions.

First, the band energies do not uniquely determine the matrix representation.  A
momentum-dependent unitary transformation can preserve all eigenvalues while changing
the eigenvectors and the Fourier coefficients.  Locality is therefore a property of a
represented Hamiltonian in a chosen gauge, not a property of the spectrum alone.
Second, removing real-space terms changes both the matrix field and its eigenvalues, but
these two errors need not rank candidate models in the same way.  Third, a small global
spectral error does not guarantee preservation of derivatives, extrema, degeneracies,
or effective masses.

This report organizes one development campaign around those distinctions.  The source
is the authenticated \(C6\)-B20-W21 rank-10 Wannier Hamiltonian for bulk silicon.  We
ask four questions:

\begin{enumerate}
  \item Can a complete represented spectrum remain unchanged while real-space locality
    and restricted frame alignment change?
  \item How much native Wigner--Seitz support is required when shells are literally
    deleted?
  \item Does bounded refitting of retained shells create useful spectral--operator
    tradeoffs?
  \item Do models selected from the 408-point development inventory remain faithful on
    dense reciprocal meshes and at conduction and valence band edges?
\end{enumerate}

The answer to the first question is yes, by construction.  The second and third answers
are more restrictive: literal deletion retains most of the native support at tight
errors, while refitting helps mainly at moderate post-hoc tolerances.  The fourth answer
is the central practical result.  Global fitting errors systematically understate
broader interpolation and band-edge consequences, and a nearly full-support model is
needed to preserve the source across all evaluated channels.

\section{Background}

\subsection{Wannier representation and gauge freedom}

For a rank-\(N\) periodic Hamiltonian, the represented matrix field is
\begin{equation}
  H(\kvec)=\sum_{\Rvec} e^{i\kvec\cdot\Rvec}T(\Rvec),
  \label{eq:fourier}
\end{equation}
where \(T(\Rvec)\) is an \(N\times N\) hopping block.  A momentum-dependent unitary
\(U(\kvec)\) defines another representation
\begin{equation}
  \widetilde H(\kvec)=U(\kvec)^\dagger H(\kvec)U(\kvec).
  \label{eq:gauge}
\end{equation}
The two fields have the same complete spectrum at each momentum, but their Fourier
blocks generally differ.  Consequently, a statement that a Hamiltonian is ``local''
requires a declared basis, gauge, lattice convention, and Fourier route.

This report uses the word \emph{gauge} only for unitary transformations inside the fixed
rank-10 represented space.  The constructed attack is not presented as a typical
Wannier90 gauge change.  It is an adversarial control that demonstrates what the
spectrum alone cannot determine.

\subsection{Matrix fidelity and spectral fidelity}

Let \(\Hsrc(\kvec)\) be the full source and \(\Hred(\kvec)\) a reduced model in the
same Wannier basis.  On the uniform 408-point development inventory, the normalized
ordered-spectrum and operator RMS losses are
\begin{align}
  \tau_1
  &=\left[\frac{1}{4080}\sum_{\kvec,n}
    \left(\frac{\epsilon_{n\kvec}^{\mathrm r}
    -\epsilon_{n\kvec}^{\mathrm W}}{\SI{1}{eV}}\right)^2\right]^{1/2},
  \label{eq:tau1}\\
  \tau_2
  &=\left[\frac{1}{4080}\sum_{\kvec}
    \frac{\|\Hred(\kvec)-\Hsrc(\kvec)\|_{\mathrm F}^2}
    {(\SI{1}{eV})^2}\right]^{1/2}.
  \label{eq:tau2}
\end{align}
The \(\SI{1}{eV}\) factor is a reporting normalization, not an acceptance threshold.
The operator loss uses identity alignment because all shell models remain in the same
frozen Wannier basis.  A pointwise fitted alignment would remove part of the model error
that the reduction is intended to measure.

Spectral error is bounded by matrix error in standard perturbation results, but the
converse does not hold.  Eigenvalues can agree while the matrix representation, frame,
and real-space coefficients differ.  The two losses are therefore retained as separate
coordinates rather than scalarized into one acceptance score.

\subsection{Why band-edge fidelity is separate}

A global RMS average gives equal formal status to all sampled bands and points.  A band
edge instead depends on a local extremum and its derivatives.  For a nondegenerate band,
the inverse effective-mass tensor is proportional to the Cartesian Hessian,
\begin{equation}
  (m^{-1})_{ij}=\frac{1}{\hbar^2}
  \frac{\partial^2 E}{\partial k_i\partial k_j}.
  \label{eq:mass}
\end{equation}
Small energy errors can produce much larger curvature errors after differentiation.

The valence edge introduces an additional difficulty.  Near a degenerate manifold,
independently ordered eigenvalues need not define analytic branches.  A basis- and
ordering-independent analysis must operate on the coupled subspace, as in multiband
\(\kvec\cdot\mathbf p\) reasoning~\cite{LuttingerKohn1955}.  Because the present source
is non-SOC, we do not assign physical heavy-hole, light-hole, or split-off labels.
Instead, we fit a matrix-valued quadratic Hamiltonian for the three-band valence
cluster.

\section{Methodology}

\subsection{Represented source and scope}

The source is the full, untruncated \(C6\)-B20-W21 Hamiltonian with rank 10 and 383
unique effective native Wigner--Seitz representatives.  It was constructed from a
Gamma-origin \(6\times6\times6\) Wannier mesh with 20 parent bands and a rank-10 target
space.  Structural checks established pair closure, real-space Hermiticity, source
reconstruction, and reciprocal periodicity.

The source is suitable for an internal reduction study but not for a claim of accurate
silicon band-edge physics.  In a separate revealed 408-point DFT-parent comparison it
missed the declared indirect-gap and effective-mass targets.  We therefore treat
\(\Hsrc\) as the represented reference and never reinterpret source-relative agreement
as agreement with DFT or experiment.

Two Fourier routes appear in the campaign and are kept distinct:
\begin{itemize}
  \item the native Wannier90 HR/Wigner--Seitz representation used for shell reduction;
  \item a centered finite-cyclic transform of sampled matrix fields, followed by
    explicit cyclic-to-ordinary conversion, used for the gauge experiment.
\end{itemize}
Native and finite-cyclic blocks are not silently substituted for one another.

\subsection{Constructed gauge attack}

The attack is a frozen rank-10, depth-2 alternating Givens circuit.  Its harmonics use
the three Cartesian-axis bonds and an axis-maximum normalization with coefficient
\(\pi/8\).  Four comparison channels are retained:
\begin{enumerate}
  \item identity coordinates;
  \item one momentum-independent global Procrustes unitary;
  \item pointwise polar alignment;
  \item the known generating inverse, used only as an oracle control.
\end{enumerate}

The reference and attacked fields are sampled on the centered \(6^3\) construction
inventory, transformed to finite-cyclic blocks, and converted to 125 ordinary
representatives with components in \(\{-2,-1,0,1,2\}\).  Complete physical-radius
supports are then evaluated on the ordered 408-point development inventory.  The study
separates complete-field gauge equivalence, conversion/interpolation error, deliberate
support truncation, and total finite-model error.

\subsection{Literal native-shell reductions}

The native representatives are grouped into 18 complete physical-radius shells.
H-shell-00 is the onsite-only control; H-shell-17 is the full 383-representative source.
The 17 proper reductions H-shell-00 through H-shell-16 are evaluated without fitting.
For a post-hoc pair \((\tau_1,\tau_2)\), the witness surface reports the smallest proper
support satisfying both inequalities.  Full support is retained as a zero-loss control
but excluded from reduced-family feasibility.

Because this family is nested and has only one support coordinate, both losses are
driven by the same omitted matrix field.  A nearly one-dimensional sequence of
requirement points is therefore expected and should not be interpreted as a general
theorem that spectral and operator losses are interchangeable.

\subsection{Refitted shell models}

To test whether the literal result is caused by the absence of adjustable parameters,
each retained complete shell is assigned one real scale,
\begin{equation}
  T_{\boldsymbol\alpha}(\Rvec)=\alpha_s T_{\mathrm W}(\Rvec),
  \qquad 0\leq\alpha_s\leq2.
  \label{eq:scaling}
\end{equation}
The same scale applies to every representative in shell \(s\), including the Hermitian
partner at \(-\Rvec\).  Blocks outside the support remain zero.  This model preserves
the fixed rank-10 Wannier basis and the within-shell block pattern; it is not a
Slater--Koster model and carries no atomic-orbital semantics.

For each proper support, 17 spectral--operator scalarization weights and three
deterministic starts are optimized with bounded L-BFGS-B and analytic gradients.  The
search produces constructive witnesses only.  A failed or absent optimization record
is not an infeasibility certificate.

\subsection{Frozen fidelity ladder}

Three converged witnesses are frozen before broader grids are inspected:
H-shell-05 with 79 representatives, H-shell-10 with 201, and H-shell-16 with 371.
No parameters are refitted later.

The exact normalized-Brillouin-zone operator RMS is obtained from the unique Fourier
coefficients by Parseval's identity.  Spectral errors are then evaluated on independent
Gamma-origin \(12^3\), \(20^3\), and \(32^3\) meshes.  Finally, the six conduction
valleys are located from sector searches and local optimization, and source-relative
gaps, positions, Hessians, and masses are calculated with identical algorithms.

\subsection{Comprehensive band-edge analysis}

For each conduction valley, Cartesian Hessians are evaluated with central-difference
steps of \(0.001\), \(0.002\), \(0.005\), and \(0.01\)~\AA$^{-1}$.  The
\(0.002\)~\AA$^{-1}$ result is primary.  In addition to full tensors and principal
masses, we report the one-valley density-of-states mass convention, the one-valley
conductivity-mass convention, a six-valley DOS-mass convention, and an equal-population
average inverse-mass convention.  These are curvature summaries, not transport or
finite-temperature carrier calculations.

At Gamma, bands 1--3 in zero-based ordering define the upper valence cluster.  We first
retain its splitting, projector distance, and principal angles.  Ordered directional
fits along [100], [110], and [111] are used only to demonstrate why individual branch
masses are unstable.

The primary valence analysis polar-transports the displaced three-band eigenspace to its
Gamma cluster and fits the even represented field to
\begin{align}
  H_{\mathrm{eff}}(\qvec)=H_{\mathrm{eff}}(0)
  &+q_x^2B_{xx}+q_y^2B_{yy}+q_z^2B_{zz} \\
  &+2q_xq_yB_{xy}+2q_xq_zB_{xz}+2q_yq_zB_{yz}.
  \label{eq:valence}
\end{align}
Thirteen unoriented rays and four positive radii per ray are used at each of three fit
windows.  Directional curvature channels are the eigenvalues of twice the contracted
quadratic matrix.  Candidate coefficient fields are polar-aligned to the source Gamma
cluster before their joint Frobenius error is calculated.

\section{Results}

\subsection{Spectrum-preserving gauge changes redistribute locality}

The maximum complete-spectrum difference between the source and attacked represented
fields is \(3.197\times10^{-14}\)~eV.  The known-inverse and pointwise-polar controls
recover the frame to near machine precision, with maximum defects of approximately
\(5.06\times10^{-16}\) and \(1.29\times10^{-15}\).  Identity and one-global-unitary
alignment leave order-one maximum frame defects of approximately 1.021 and 1.020.

\resultfigure
  {../../analysis/c6-b20-w21-gauge-alignment-locality/figures/alignment-invariance.pdf}
  {Complete spectral invariance and alignment diagnostics for the frozen gauge attack.}
  {Pointwise polar and known-inverse channels are oracle controls; identity and one
  global unitary cannot remove the momentum-dependent attack.}
  {fig:alignment}

Despite complete spectral invariance, the non-onsite finite-block power changes from
\(20.7574\) to \(21.5057\)~eV\(^2\) per retained state.  The shell-resolved and cumulative
power distributions change accordingly.  This is a direct counterexample to the idea
that the complete spectrum fixes real-space locality.

\resultfigure
  {../../analysis/c6-b20-w21-gauge-alignment-locality/figures/fourier-power-redistribution.pdf}
  {Redistribution of finite Fourier-block power under the spectrum-preserving attack.}
  {The cyclic-to-ordinary removed power is shown separately and is not counted as
  deliberate support truncation.}
  {fig:power}

Finite-support requirements also change.  The reference and attacked support-error
curves have different non-dominated supports even though their complete represented
spectra agree.  The result concerns the constructed attack, not the distribution of
ordinary Wannier90 gauges.

\resultfigure
  {../../analysis/c6-b20-w21-gauge-alignment-locality/figures/pareto-support-error-frontiers.pdf}
  {Discrete complete-support error frontiers for reference and attacked representations.}
  {Markers identify observed non-dominated supports; connecting segments are visual
  guides and do not interpolate unobserved support choices.}
  {fig:gauge-frontier}

\subsection{Literal shell deletion retains most of the source at tight errors}

Among all proper literal reductions, H-shell-16 minimizes both primary RMS losses while
retaining 371 of 383 representatives:
\begin{equation}
  (\tau_1,\tau_2)=
  (4.4636\times10^{-4},1.5173\times10^{-3}).
\end{equation}
Thus the literal nested family does not display a global spectral--operator conflict at
its tightest proper reduction.  Smaller supports show local nonmonotonicity, but the
requirement points remain close to a line because one nested support coordinate controls
both errors.

The post-hoc pair \((0.1,0.15)\) first admits 135 representatives; \((0.01,0.02)\)
requires 313; and \((0.001,0.002)\) requires 371.  No proper literal reduction satisfies
\((0.0004,0.001)\).  These coordinates summarize the observed family and are not
prospective acceptance thresholds.

\resultfigure
  {../../analysis/c6-b20-w21-tau1-tau2-posthoc-exploration/figures-v3/single-witness-surface.pdf}
  {Minimum-support witness surface for literal complete-shell deletion.}
  {The full source is excluded from reduced-family feasibility.  Gray cells contain no
  proper literal witness at the displayed post-hoc pair.}
  {fig:literal-surface}

\subsection{Shell refitting creates genuine two-loss tradeoffs}

Bounded shell scaling changes the geometry of the candidate family.  A fixed support can
now trace a spectral--operator tradeoff arc rather than contribute only one deletion
point.  Across 444 converged witnesses, refitting reduces the minimum observed support
at several illustrative pairs (Table~\ref{tab:witnesses}).

\begin{table}[ht]
  \centering
  \caption{Observed minimum-support literal and refitted witnesses.  Threshold pairs are
  post-hoc descriptive coordinates, not acceptance limits.}
  \label{tab:witnesses}
  \begin{tabular}{@{}lrr@{}}
    \toprule
    $(\tau_1,\tau_2)$ & Literal representatives & Refitted representatives \\
    \midrule
    $(1.0,1.0)$       & 19  & 13  \\
    $(0.1,0.15)$      & 135 & 79  \\
    $(0.05,0.1)$      & 177 & 135 \\
    $(0.02,0.04)$     & 225 & 201 \\
    $(0.01,0.02)$     & 313 & 313 \\
    $(0.001,0.002)$   & 371 & 371 \\
    \bottomrule
  \end{tabular}
\end{table}

Within H-shell-16, an operator-oriented witness attains approximately
\((4.04\times10^{-4},1.394\times10^{-3})\), while a spectral-oriented witness attains
\((2.98\times10^{-4},1.745\times10^{-3})\).  This is a genuine same-support tradeoff,
although bounded multistart optimization does not establish global optimality.

\resultfigure
  {../../analysis/c6-b20-w21-refitted-ws-shell-witness-posthoc/figures-v2/single-refitted-witness-surface.pdf}
  {Minimum-support surface over converged refitted shell witnesses.}
  {Gray means that the bounded search found no converged witness, not that the pair is
  mathematically infeasible.}
  {fig:refitted-surface}

\subsection{The 408-point inventory understates broader error}

The frozen fidelity ladder gives the central quantitative comparison.  For every model,
the fitted-408 spectral and operator losses are smaller than their broader counterparts
(Table~\ref{tab:fidelity}).

\begin{table}[ht]
  \centering
  \caption{Fitted-inventory and broader frozen-model RMS errors.}
  \label{tab:fidelity}
  \small
  \begin{tabular}{@{}lrrrrrr@{}}
    \toprule
    & \multicolumn{3}{c}{Spectral RMS (eV)}
    & \multicolumn{3}{c}{Operator RMS (eV)} \\
    \cmidrule(lr){2-4}\cmidrule(l){5-7}
    Model & 408 & $32^3$ & Ratio & 408 & Parseval & Ratio \\
    \midrule
    H-shell-05 & 0.07242  & 0.18196  & 2.51 & 0.12814  & 0.26004  & 2.03 \\
    H-shell-10 & 0.01320  & 0.02785  & 2.11 & 0.03229  & 0.05956  & 1.84 \\
    H-shell-16 & 0.000404 & 0.000657 & 1.62 & 0.001394 & 0.001747 & 1.25 \\
    \bottomrule
  \end{tabular}
\end{table}

The sampled operator RMS agrees with the exact Parseval value to at most
\(5.56\times10^{-17}\)~eV.  Spectral RMS changes between the \(12^3\) and \(32^3\)
meshes remain below about 1.4\%, so the dense values are stable over the tested mesh
sequence.  The discrepancy is therefore not a late-grid fluctuation; it is a difference
between fitting-inventory and broader-domain evaluation.

\resultfigure
  {../../analysis/c6-b20-w21-wannier-only-fidelity-ladder/figures/wannier-fidelity-ladder.pdf}
  {Frozen-model fidelity from the fitted inventory to exact and dense evaluations.}
  {Supports and shell scales remain unchanged at every evaluation level.}
  {fig:fidelity-ladder}

\subsection{Band-edge features are more demanding than global RMS}

H-shell-05 changes the source-relative indirect gap by \(-65.8\)~meV, displaces a
conduction valley by as much as \(0.0948\)~\AA$^{-1}$, and changes at least one principal
mass by 67.7\%.  H-shell-10 has a \(32^3\) all-band spectral RMS below 0.028~eV but still
changes a principal mass by 53.2\%.  H-shell-16 changes the gap by only \(-0.503\)~meV,
the maximum valley position by \(0.00276\)~\AA$^{-1}$, and the maximum principal mass by
2.54\%.

\resultfigure
  {../../analysis/c6-b20-w21-wannier-only-fidelity-ladder/figures/source-relative-feature-fidelity.pdf}
  {Source-relative conduction-edge consequences of the frozen shell models.}
  {The reference is the full Wannier source, not DFT or experiment.}
  {fig:feature-fidelity}

The expanded tensor analysis leads to the same ordering.  The source six-valley
DOS-mass convention is \(1.4764\,m_e\), and its equal-population conductivity-mass
convention is \(0.25063\,m_e\).  H-shell-05 changes them by 16.9\% and 20.4\%; H-shell-10
by 22.9\% and 6.90\%; and H-shell-16 by 0.333\% and 0.0216\%.  The maximum mass variation
across four finite-difference steps is below 0.172\% for every model, so stencil choice
does not explain the large moderate-support errors.

\resultfigure
  {../../analysis/c6-b20-w21-wannier-only-comprehensive-band-edge/figures/derived-electron-mass-fidelity.pdf}
  {Derived conduction-valley mass conventions and finite-difference sensitivity.}
  {The DOS and conductivity quantities are curvature conventions, not carrier or
  transport calculations.}
  {fig:electron-masses}

\subsection{Moderate reduction damages the degenerate valence edge}

The full source has a three-band Gamma splitting of only 0.02495~meV.  H-shell-05 and
H-shell-10 increase this to 102.2 and 54.8~meV.  H-shell-16 reduces the error
substantially but still yields 2.50~meV.  The corresponding maximum source-cluster
principal angles are 1.856, 0.925, and 0.0420 degrees.  The moderate models therefore
preserve a recognizable three-dimensional subspace while strongly changing the
Hamiltonian inside it.

\resultfigure
  {../../analysis/c6-b20-w21-wannier-only-comprehensive-band-edge/figures/gamma-valence-cluster-fidelity.pdf}
  {Gamma valence-cluster splitting and subspace rotation.}
  {Subspace agreement and preservation of internal degeneracy are separate conditions.}
  {fig:gamma-cluster}

Independently ordered source branches do not yield stable hole masses: the upper branch
is not hole-like for every tested window in any of the three directions, and finite
ordered-branch masses vary by up to about 75.5\%.  The primary multiband fit avoids that
assignment.  At the 0.01~\AA$^{-1}$ window, its source curvature-channel masses are
listed in Table~\ref{tab:valence}.

\begin{table}[ht]
  \centering
  \caption{Full-source non-SOC valence curvature-channel masses in electron-mass units.}
  \label{tab:valence}
  \begin{tabular}{@{}lrrr@{}}
    \toprule
    Direction & Channel 1 & Channel 2 & Channel 3 \\
    \midrule
    $[100]$ & 0.1411 & 0.2771 & 0.3015 \\
    $[110]$ & 0.09690 & 0.3119 & 2.923 \\
    $[111]$ & 0.08770 & 0.6971 & 0.7722 \\
    \bottomrule
  \end{tabular}
\end{table}

These channels belong to a scalar-relativistic three-band quadratic operator.  They are
not physical SOC heavy-hole, light-hole, or split-off masses.  After cluster alignment,
H-shell-05, H-shell-10, and H-shell-16 have quadratic-coefficient errors of 9.54\%,
5.88\%, and 0.279\%, and maximum curvature-channel mass changes of 38.9\%, 19.7\%, and
4.53\%.

\resultfigure
  {../../analysis/c6-b20-w21-wannier-only-comprehensive-band-edge/figures/multiband-valence-fidelity-summary.pdf}
  {Summary of non-SOC multiband valence fidelity.}
  {Gamma splitting, aligned quadratic-operator error, and maximum curvature-channel
  change expose distinct failure modes.}
  {fig:valence-summary}

A related symmetry diagnostic reaches the same conclusion.  The full source's direct-gap
spread among equivalent X points is 0.00366~meV and among L points is 0.0239~meV.
H-shell-05 increases those spreads to 150.1 and 68.7~meV; H-shell-10 to 36.7 and
63.9~meV; and H-shell-16 to 0.427 and 2.97~meV.  Moderate shell refitting can therefore
amplify small residual symmetry differences through cancellations among retained
blocks.

\resultfigure
  {../../analysis/c6-b20-w21-wannier-only-comprehensive-band-edge/figures/high-symmetry-edge-fidelity.pdf}
  {Direct-gap and equivalent-point symmetry diagnostics.}
  {The displayed spreads characterize the represented models and are not an independent
  crystal-symmetry validation.}
  {fig:symmetry}

\clearpage
\section{Discussion}

\subsection{Locality is representation-dependent}

The gauge experiment supplies the cleanest conceptual result.  The complete spectrum is
unchanged to roundoff, yet block power and finite-support requirements change.  Thus a
band structure cannot determine a unique hopping-locality profile.  Any comparison of
``short-range'' Hamiltonians must identify the represented state space and the allowed
gauge transformations before locality is interpreted.

This does not imply that locality is arbitrary or useless.  A fixed construction,
gauge, and basis produce a well-defined optimization problem.  It does imply that
locality claims cannot be elevated from spectral data alone, and that a constructed
gauge counterexample should not be confused with ordinary variation between independently
localized Wannier calculations.

\subsection{Refitting helps, but support remains the dominant resource}

Literal shell deletion is too restrictive to reveal a strong global two-loss conflict:
one nested coordinate controls both matrix and spectral losses.  Adding bounded shell
scales exposes genuine tradeoff arcs and reduces support at moderate post-hoc error
pairs.  This shows that some error from deletion can be compensated by renormalizing
retained shells.

The compensation has limits.  At tight errors the observed minimum support remains 371
representatives, close to the 383-representative source.  Moreover, the bounded search
is constructive rather than exhaustive.  The correct conclusion is that refitting
produces useful witnesses, not that the displayed surface is a global optimum or that
gray cells are infeasible.

\subsection{A fitting inventory is not an interpolation guarantee}

Every frozen model performs worse on the exact integrated or independent dense
evaluation than on the 408-point inventory used to characterize fitted witnesses.  The
discrepancy is largest at moderate support and decreases near full support.  This is the
expected role of a broader evaluation: it reveals behavior between fitted points without
changing the model.

No held-out set is required to make the finite statement that a model attains a given
loss on the 408 retained points.  A held-out or all-\(\kvec\) argument is required only
when the claim expands to interpolation or Brillouin-zone behavior.  Parseval's identity
provides the exact integrated operator statement here, while dense meshes provide the
spectral interpolation check.

\subsection{Band edges should be treated as first-class objectives}

The conduction results show why all-band RMS is an incomplete reduction target.
H-shell-10 appears moderately accurate in the dense global spectral metric, yet changes
a principal electron mass by more than 50\%.  Effective masses depend on second
derivatives and therefore amplify local shape errors that contribute little to a global
average.

The valence result is stronger.  Independent band Hessians are not merely noisy; at the
near-degenerate source edge they are not a stable mathematical object.  The appropriate
quantity is a matrix-valued multiband quadratic Hamiltonian.  Moderate reductions can
make individual branches look numerically stable by incorrectly splitting the
manifold.  Stability after symmetry breaking is not fidelity.

For this non-SOC source, the multiband operator is the strongest justified hole-sector
result.  Physical silicon heavy-hole, light-hole, split-off, or acceptor claims would
require an explicitly spinor, SOC-enabled source and a separately qualified comparison.
SOC is not required for the internal compression conclusion reached here.

\subsection{What the study establishes}

The retained calculations support the following bounded conclusions:
\begin{enumerate}
  \item Complete spectral equivalence does not determine real-space block locality.
  \item Operator and spectral losses are distinct and should remain separate reduction
    coordinates.
  \item Shell refitting can reduce support at moderate post-hoc losses, but tight
    internal fidelity still requires near-full native support.
  \item Errors on the fitted 408-point inventory understate broader interpolation error.
  \item Global spectral RMS does not control conduction masses, valence degeneracy, or
    multiband curvature.
  \item Among the frozen witnesses, H-shell-16 is consistently closest to the source,
    while H-shell-05 and H-shell-10 are not adequate band-edge reductions.
\end{enumerate}

\subsection{Limitations and next steps}

The principal limitation is the represented source itself.  It fails separate
DFT-parent gap and electron-mass criteria, so source-relative preservation does not
validate silicon physics.  The source is also scalar-relativistic and non-SOC.  The
analyses are producer-checked but not independently reconstructed, and all tolerance
surfaces are post hoc rather than prospectively accepted.

A scientifically stronger continuation would keep the current models frozen and add,
in order:
\begin{enumerate}
  \item independent reconstruction of the retained analyses;
  \item a clearly separated overlay against the existing revealed 408-point DFT parent,
    without refitting to that overlay;
  \item a qualified SOC DFT/Wannier branch if physical hole or acceptor properties are
    required; and
  \item prospective reduction thresholds declared before inspecting any future
    confirmation set.
\end{enumerate}
None of these steps changes the present internal conclusion: for the tested source and
model family, near-full support is required to preserve matrix, spectral, symmetry, and
band-edge behavior simultaneously.

\subsection{Reproducibility}

Canonical requests, authenticated source identities, producer programs, machine-readable
results, figure manifests, and execution records are retained under
\path{calculations/Si/bulk/analysis/}.  The gauge, literal-shell, refitted-shell,
fidelity-ladder, and comprehensive-band-edge analyses each retain their own provenance
boundary.  No numerical payload in this report is the sole record of a result.

\bibliographystyle{unsrt}
\bibliography{references}

\end{document}
